\documentclass[11pt,a4paper]{article}

% ===== XeLaTeX + Korean =====
\usepackage{kotex}
\usepackage{fontspec}
\setmainfont{Times New Roman}
\setmainhangulfont{AppleMyungjo}
\setsanshangulfont{Apple SD Gothic Neo}

% ===== Layout =====
\usepackage[a4paper,top=18mm,bottom=18mm,left=20mm,right=20mm]{geometry}
\usepackage{fancyhdr}
\setlength{\headheight}{24pt}
\usepackage{enumitem}
\usepackage{mdframed}
\usepackage{array}
\usepackage{booktabs}

\setlength{\parindent}{1em}
\linespread{1.18}

% ===== Header / Footer =====
\pagestyle{fancy}
\fancyhf{}
\renewcommand{\headrulewidth}{0.6pt}
\fancyhead[L]{\sffamily\bfseries 2026 고1 영어 변형 — 정답 및 해설 지문 34}
\fancyhead[R]{\sffamily 결과에 따른 행동 조정}
\renewcommand{\footrulewidth}{0pt}
\fancyfoot[C]{\framebox[16mm][c]{\thepage}}

% ===== helpers =====
\newcommand{\qno}[1]{\par\vspace{10pt}\noindent{\sffamily\bfseries #1.}\ }
\newmdenv[linewidth=0.6pt,roundcorner=3pt,innertopmargin=7pt,innerbottommargin=7pt,
          innerleftmargin=9pt,innerrightmargin=9pt,skipabove=5pt,skipbelow=5pt]{pbox}

\begin{document}

\begin{center}
{\fontsize{15}{17}\selectfont\sffamily\bfseries [지문 34] 정답 및 해설}
\end{center}
\vspace{6pt}

% ===== 정답 표 =====
\begin{center}
\renewcommand{\arraystretch}{1.4}
\begin{tabular}{ccccc}
\toprule
문항 & 유형 & 정답 & 배점 \\
\midrule
1 & 빈칸추론 & ③ & 3점 \\
2 & 영어 주제 & ① & 2점 \\
3 & 서술형 — 해석 & 모범답안 참조 & 2점 \\
4 & 내용 불일치 & ④ & 2점 \\
\bottomrule
\end{tabular}
\end{center}

\vspace{8pt}

% ===== Q1 해설 =====
\qno{1} \textsf{\bfseries 빈칸추론 — 정답: ③ advantage}

\begin{pbox}
\textbf{정답 근거}: ``we tend to do things that bring us some kind of \underline{benefit}'' 가 원문 표현이며, 빈칸은 이 ``benefit(이익)''을 대체한다. 바로 이어지는 문장에서 ``help us avoid some kind of problem, discomfort, or disadvantage''와 대조를 이루므로, 긍정적 결과를 나타내는 ③ \textbf{advantage}(이점, 유리함)가 정답이다.

\medskip
\textbf{오답 소거}:
\begin{itemize}[leftmargin=1.5em,itemsep=1pt,topsep=2pt]
  \item ① difficulty — 어려움은 우리가 추구하는 것이 아님. 오답.
  \item ② uncertainty — 불확실성은 회피 대상. 오답.
  \item ④ obligation — 의무감에 의한 행동은 본문의 ``choosing to do'' 맥락과 다름. 매력적 오답이나, 본문은 자발적 선택을 설명함.
  \item ⑤ distraction — 주의 분산은 부정적 개념. 오답.
\end{itemize}
\end{pbox}

% ===== Q2 해설 =====
\qno{2} \textsf{\bfseries 영어 주제 — 정답: ①}

\begin{pbox}
\textbf{정답 근거}: 지문 전체는 인간이 결과(consequences)에 따라 행동을 조정한다는 내용이며, ① ``how the consequences of behavior shape what people choose to do''가 이를 가장 정확하게 요약한다.

\medskip
\textbf{오답 소거}:
\begin{itemize}[leftmargin=1.5em,itemsep=1pt,topsep=2pt]
  \item ② 아동 칭찬 사례는 하나의 예시일 뿐, 글 전체의 주제가 아님.
  \item ③ 동기와 훈련의 차이는 본문에서 다루지 않음.
  \item ④ 불편함이 성취로 이어진다는 내용은 본문과 정반대 (본문은 불편함을 \textit{회피}한다고 서술).
  \item ⑤ 반복의 역할은 댄서 예시에 국한되며 주제가 아님.
\end{itemize}
\end{pbox}

% ===== Q3 해설 (서술형) =====
\qno{3} \textsf{\bfseries 서술형 — 해석}

\begin{pbox}
\textbf{대상 문장}:\par
\hspace{1em}\textit{A dancer who feels that a particular conditioning exercise is not giving any results is likely to give that exercise up in favor of another one that seems more promising.}

\medskip
\textbf{모범 답안}:\par
\hspace{1em}특정 컨디셔닝 운동이 아무 결과도 내지 못하고 있다고 느끼는 댄서는, 더 가능성 있어 보이는 다른 운동을 위해 그 운동을 포기할 것 같다.

\medskip
\textbf{채점 포인트} (각 1점씩, 총 2점):
\begin{itemize}[leftmargin=1.5em,itemsep=1pt,topsep=2pt]
  \item \textbf{[1점]} ``give that exercise up in favor of another one'' = ``다른 운동을 위해 그 운동을 포기하다'' (in favor of 의미 정확히 반영)
  \item \textbf{[1점]} ``that seems more promising'' = ``더 가능성 있어 보이는'' 또는 ``더 유망해 보이는'' (관계절 구조 및 의미 정확히 반영)
\end{itemize}
\end{pbox}

\vspace{8pt}
\qno{4} \textsf{\bfseries 내용 불일치 - 정답: ④}
\begin{pbox}
본문은 사람들이 이익을 주거나 문제와 불편을 피하게 해 주는 행동을 더 많이 하고, 보상이나 이익을 주지 않는 행동은 피한다고 말한다. 따라서 보상 없이 불편을 만드는 방식으로 행동하는 경향이 있다는 ④는 본문과 반대이다.
\end{pbox}

\end{document}
